% =====================================================================
%  02 模型假设
%  写出建模时做的合理假设,每条假设在建模中会被用到,
%  尽量每条都能对应到 04 模型建立里的某个式子/符号。
%
%  符号说明表:由队长(Claude)统一维护,全队以此为准。
%  建模手写公式、代码手写代码、论文手排版都用同一套符号。
%  新符号出现时,更新此表并推送,大家在群里同步。
% =====================================================================
\section{模型假设}

\begin{enumerate}
  \item 假设一:……
  \item 假设二:……
  \item 假设三:……
\end{enumerate}

\newpage
\section*{符号说明}

\begin{table}[htbp]
\centering
\caption{主要符号说明}
\begin{tabular}{lll}
\toprule
符号 & 含义 & 单位 \\
\midrule
$\alpha$   & 针肋宽度比(针肋直径/通道宽度) & — \\
$\beta$    & 歧管深高比(歧管层/微通道层深高比) & — \\
$n$        & 单个歧管单元内沿流向的针肋排数 & — \\
$R_{th}$   & 总热阻 & K/W \\
$\Delta P$ & 压降 & Pa \\
$\sigma_T$ & 温度非均匀性指标 & K \\
$T_{max}$  & 芯片最高温度 & K \\
$T_{in}$   & 冷却液入口温度 & K \\
$Q$        & 热流密度/总热功率 & W/m$^2$ 或 W \\
$h$        & 对流换热系数 & W/(m$^2$$\cdot$K) \\
$A$        & 换热面积 & m$^2$ \\
$D_h$      & 通道水力直径 & m \\
$u$        & 流速 & m/s \\
$\rho$     & 冷却液密度 & kg/m$^3$ \\
$c_p$      & 冷却液比热容 & J/(kg$\cdot$K) \\
$\mu$      & 冷却液动力黏度 & Pa$\cdot$s \\
$f$        & 沿程阻力系数 & — \\
$Re$       & 雷诺数 & — \\
\bottomrule
\end{tabular}
\end{table}

% 说明:以上为预填常用符号,比赛过程中随建模进度增删。
% 新增符号时:告知队长 → 队长更新本表 → 推送,全队同步。
